\documentclass[11pt]{article}
\usepackage[T1]{fontenc}
\usepackage{amsmath,amssymb}
\usepackage[margin=1in]{geometry}
\title{Mathematics in LaTeX: Beginner's Guide}
\author{}
\date{}
\begin{document}
\maketitle
\tableofcontents
\bigskip

\section{Inline and Display Mathematics}
\textbf{Definition.} LaTeX has two common ways to write mathematical expressions. Inline mathematics sits within a sentence. Display mathematics appears separately, usually centered on its own line.

\textbf{Example 1: Inline mathematics}
\begin{verbatim}
The equation $a+b=c$ is an example.
\end{verbatim}
\textbf{Expected result:} The equation $a+b=c$ is an example.

\textbf{Example 2: Display mathematics}
\begin{verbatim}
\[
a+b=c
\]
\end{verbatim}
\textbf{Expected result:}
\[
a+b=c
\]
\textbf{Commands.} Dollar signs \verb|$...$| mark inline maths. The commands \verb|\[| and \verb|\]| begin and end an unnumbered display equation.

\textbf{Practice.} Write $x+y=z$ inside a sentence and then on its own line.

\section{Superscripts and Subscripts}
\textbf{Definition.} A superscript is placed above a symbol, often to show a power. A subscript is placed below a symbol, often to label it.
\textbf{Source code:}
\begin{verbatim}
\[ x^2 \qquad a_1 \qquad x_1^2 \]
\[ x^{10} \qquad a_{12} \]
\end{verbatim}
\textbf{Expected result:}
\[ x^2 \qquad a_1 \qquad x_1^2 \]
\[ x^{10} \qquad a_{12} \]
\textbf{Commands.} The caret \verb|^| creates a superscript; the underscore \verb|_| creates a subscript. Use braces to group more than one character, as in \verb|x^{10}|.
\textbf{Practice.} Write $y^3$, $b_2$, and $m_{15}$.

\section{Fractions and Roots}
\textbf{Definition.} A fraction has a numerator (top part) and denominator (bottom part). A root shows a value that, when raised to a power, gives the original number.
\textbf{Source code:}
\begin{verbatim}
\[ \frac{a}{b} \qquad \sqrt{x} \qquad \sqrt[3]{x} \]
\end{verbatim}
\textbf{Expected result:}
\[ \frac{a}{b} \qquad \sqrt{x} \qquad \sqrt[3]{x} \]
\textbf{Commands.} \verb|\frac{a}{b}| creates a fraction. \verb|\sqrt{x}| creates a square root. The optional index in \verb|\sqrt[3]{x}| makes a cube root.
\textbf{Practice.} Write one-half, the square root of $a+b$, and the fourth root of $x$.

\section{Greek Letters}
\textbf{Definition.} Greek letters are commonly used for mathematical quantities, angles, and constants.
\textbf{Source code:}
\begin{verbatim}
\[ \alpha,\ \beta,\ \gamma,\ \delta \]
\[ \pi,\ \theta,\ \lambda,\ \omega \]
\end{verbatim}
\textbf{Expected result:}
\[ \alpha,\ \beta,\ \gamma,\ \delta \]
\[ \pi,\ \theta,\ \lambda,\ \omega \]
\textbf{Commands.} A backslash followed by a letter name produces the symbol: \verb|\alpha| gives $\alpha$, and \verb|\theta| gives $\theta$. The command \verb|\pi| gives $\pi$.
\textbf{Practice.} Find and write the commands for $\mu$, $\sigma$, $\phi$, and $\Delta$.

\section{Mathematical Operators}
\textbf{Definition.} Operators represent mathematical actions. Summation adds terms, product multiplies terms, and a limit describes a value an expression approaches.
\textbf{Source code:}
\begin{verbatim}
\[ \sum_{i=1}^{n} i \]
\[ \prod_{i=1}^{n} i \]
\[ \lim_{x \to 0} f(x) \]
\end{verbatim}
\textbf{Expected result:}
\[ \sum_{i=1}^{n} i \]
\[ \prod_{i=1}^{n} i \]
\[ \lim_{x \to 0} f(x) \]
\textbf{Commands.} \verb|\sum| makes a summation sign, \verb|\prod| makes a product sign, and \verb|\lim| makes a limit operator. Underscores and carets set the lower and upper limits.
\textbf{Practice.} Write a summation from $i=1$ to $5$ and a limit as $x$ approaches $2$.

\section{Brackets and Grouping}
\textbf{Definition.} Brackets group parts of an expression. LaTeX can also resize brackets to fit tall expressions such as fractions.
\textbf{Source code:}
\begin{verbatim}
\[ (a+b) \qquad [x+y] \]
\[ \left(\frac{a}{b}\right) \]
\[ \left[\frac{x+1}{x-1}\right] \]
\end{verbatim}
\textbf{Expected result:}
\[ (a+b) \qquad [x+y] \]
\[ \left(\frac{a}{b}\right) \]
\[ \left[\frac{x+1}{x-1}\right] \]
\textbf{Commands.} Type ordinary parentheses and square brackets directly. Pair \verb|\left| with \verb|\right| to resize matching delimiters around an expression.
\textbf{Practice.} Put large round brackets around a fraction and large square brackets around a sum.

\section{Equations and Alignment}
\textbf{Definition.} Displayed equations appear on separate lines. An alignment environment lines up related equations at a chosen point, commonly the equals sign.
\textbf{Source code:}
\begin{verbatim}
\begin{align*}
x + 2 &= 5 \\
x &= 5 - 2 \\
x &= 3
\end{align*}
\end{verbatim}
\textbf{Expected result:}
\begin{align*}
x + 2 &= 5 \\
x &= 5 - 2 \\
x &= 3
\end{align*}
\textbf{Commands.} The \verb|align*| environment (from the \verb|amsmath| package) aligns multiple lines without numbering them. The ampersand \verb|&| marks the alignment point, and \verb|\\| starts a new row. Add \verb|\usepackage{amsmath}| to the document preamble.
\textbf{Practice.} Show the steps for solving $y+4=9$, aligning each equals sign.

\section{Matrices}
\textbf{Definition.} A matrix is a rectangular arrangement of entries in rows and columns.
\textbf{Source code:}
\begin{verbatim}
\[
\begin{bmatrix}
1 & 2 \\
3 & 4
\end{bmatrix}
\]
\end{verbatim}
\textbf{Expected result:}
\[
\begin{bmatrix}
1 & 2 \\
3 & 4
\end{bmatrix}
\]
\textbf{Commands.} The \verb|bmatrix| environment creates a matrix with square brackets. Use \verb|&| to separate columns and \verb|\\| to start a new row. It is provided by the \verb|amsmath| package.
\textbf{Practice.} Create a matrix with two rows and three columns containing the numbers 1 to 6.
\end{document}
